% Systematic review template (LaTeX, PRISMA-style headings)
% Use this when you need a reproducible search and screening protocol.

\documentclass[11pt]{article}

\usepackage[T1]{fontenc}
\usepackage[utf8]{inputenc}
\usepackage{lmodern}
\usepackage{microtype}
\usepackage{amsmath,amssymb}
\usepackage{booktabs}
\usepackage{hyperref}
\usepackage[numbers]{natbib}

\title{Title of Systematic Review}
\author{Author Name}
\date{\today}

\begin{document}
\maketitle

\begin{abstract}
Summarize: (1) objective, (2) methods (databases, dates, criteria), (3) key findings/themes, (4) main gaps and implications.
\end{abstract}

\section{Introduction}
State the research question and why a systematic review is needed.

\section{Methods}

\subsection{Eligibility Criteria}
Define inclusion/exclusion criteria (time window, languages, study types, task scope).

\subsection{Information Sources}
List databases and any additional sources (backward/forward citation chasing).

\subsection{Search Strategy}
Provide the exact query strings and the dates executed.

\subsection{Screening and Selection}
Describe title/abstract and full-text screening, duplicates handling, and reasons-for-exclusion logging.

\subsection{Data Extraction}
Define the extraction fields (method, data, metrics, results, limitations, reproducibility artifacts).

\subsection{Quality Appraisal}
State the appraisal rubric and how disagreements were resolved.

\section{Results}
Report selection counts and summarize themes. Include an evidence table.

\section{Discussion}
Interpret findings, highlight contradictions and limitations, and propose future work.

\section{Conclusion}
Concise summary of what is known and what remains open.

\bibliographystyle{plainnat}
\bibliography{references}

\end{document}

